% mode: ini-etex
% Arguments and skipped text read from token lists in alignments.  §364 fires on a tab or a \cr at align_state 0, for a character or
% a control sequence alike, and a capture loop must not pass over one.  What each case does is in the \message before it.
\input prelude
\tracingmacros=1 \errorcontextlines=100
% These do not typeset what they took, since a tab in a body would be an error of its own: \tracingmacros shows the arguments.
\def\dm#1.{}
\def\um#1{}
\def\tm#1ab#2{}
\def\pass#1{#1}
\let\amp=&
\let\crr=\cr
\let\spn=\span
\def\row#1{\def\q{#1}\setbox0\vbox{\halign{##&##\cr\q\cr}}}
\def\case#1#2{\noindent\par\message{^^Jcase #1:}\row{#2}\showbox0 }

% At align_state 0, reached by a letter in the cell, so that the u template has ended.
\case{tab in a delimited argument}{x\dm a&b.}
\case{tab as an undelimited argument}{x\um&}
\case{tab inside braces is only a token}{x\um{a&b}}
\case{tab in a group in a delimited argument}{x\dm{a&b}.}
\case{a group and more in a delimited argument}{x\dm{a}&{b}.}
\case{a control sequence meaning tab}{x\dm a\amp b.}
\case{the same as a whole argument}{x\um\amp}
\case{and in braces}{x\um{a\amp b}}
\case{cr in a delimited argument}{x\dm a\cr b.}
\case{let to cr}{x\dm a\crr b.}
\case{cr as an argument}{x\um\cr}
\case{let to cr, as an argument}{x\um\crr}
\case{cr in braces}{x\um{\cr}}
\case{span in a delimited argument}{x\dm a\span b.}
\case{let to span}{x\dm a\spn b.}
\case{a two-token delimiter}{x\tm a&bab z}
\case{the delimiter itself is a tab}{\def\td##1&{[##1]}x\td ab&c}
\case{a space before the argument, then a tab}{x\um &}
\case{the first argument is fine, the second a tab}{\def\tw##1##2{[##1|##2]}x\tw{a}&}
\case{a macro in braces that then gets its tab}{x\um{\um}&}

% The same before the first letter: align_peek reads these at align_state 1000000.
\case{align_peek expands the macro}{\dm a&b.}
\case{align_peek, undelimited}{\um&}
\case{align_peek, cr}{\dm a\cr b.}
\case{align_peek, skipping}{\iffalse&\fi}

% align_state -1 and 1: a skipped brace moves it, then a group takes it through 0.  This text is not a macro's argument, as its braces do not match.
\def\skipr{\noindent\par\message{^^Jthe cell has gone below 0:}}
\def\ua{\um{a&b}}
\def\ub{\dm{a&b}.}
\def\uc{\um{{a&b}c&d}}
\skipr\setbox0\vbox{\halign{#&#\cr x\iffalse}\fi\ua\iffalse{\fi\cr}}\showbox0
\skipr\setbox0\vbox{\halign{#&#\cr x\iffalse}\fi\ub\iffalse{\fi\cr}}\showbox0
\skipr\setbox0\vbox{\halign{#&#\cr x\iffalse}}\fi\uc\iffalse{{\fi\cr}}\showbox0
\def\skipl{\noindent\par\message{^^Jthe cell has gone above 0:}}
\def\ud{\dm a&b.}
\def\ue{\um&}
\def\uf{\dm a\cr b.}
\skipl\setbox0\vbox{\halign{#&#\cr x\iffalse{\fi\ud\iffalse}\fi\cr}}\showbox0
\skipl\setbox0\vbox{\halign{#&#\cr x\iffalse{\fi\ue\iffalse}\fi\cr}}\showbox0
\skipl\setbox0\vbox{\halign{#&#\cr x\iffalse{\fi\uf\iffalse}\fi\cr}}\showbox0

% Skipped text.
\case{a tab in skipped text}{x\iffalse&\fi}
\case{a tab in a group in skipped text}{x\iffalse{&}\fi}
\case{a left brace then a right brace, then a tab}{x\iffalse{}&\fi}
\case{cr in skipped text}{x\iffalse\cr\fi}
\case{let to cr in skipped text}{x\iffalse\crr\fi}
\case{span in skipped text}{x\iffalse\spn\fi}
\case{an else branch skipped}{x\iftrue\else&\fi}
\case{ifcase}{x\ifcase2 a&\or b\or c\fi}
\case{a tab in an argument that is skipped}{\def\w##1{x\iffalse##1\fi}\w{&}}
\case{an argument that is skipped, in a group}{\def\w##1{x\iffalse{##1}\fi}\w{&}}
\case{an argument that is skipped outside the state of 0}{\def\w##1{\iffalse##1\fi}\w{&}}

% Skipped text with braces that do not match, from a list.  A group in front keeps the text balanced; \egroup ends it after the skip.
\case{a brace in front, right brace and tab skipped}{x{\iffalse}&\fi\egroup}
\case{a brace in front, right, left, tab}{x{\iffalse}{&\fi\egroup\iffalse}\fi}
\case{a brace in front, right, cr}{x{\iffalse}\cr\fi\egroup}
\case{right, cr by another name}{x{\iffalse}\crr\fi\egroup}
\case{a parameter holds the tab}{\def\w##1{x{\iffalse}##1\fi\egroup}\w{&}}
\case{a parameter holds the right brace}{\def\w##1{x{\iffalse##1&\fi\egroup}\w{}}
\case{an argument after a skipped right brace, undelimited}{x{\iffalse}\fi\um&\egroup}
\case{delimited}{x{\iffalse}\fi\dm a&b.\egroup}
\case{in a group that reaches 0}{x\iffalse{\fi{\iffalse}}\fi\um{a&b}\egroup\iffalse}\fi}
% And in a template.
\def\tpl#1#2{\noindent\par\message{^^Jcase #1:}\setbox0\vbox{\halign{#2\cr x\cr}}\showbox0 }
% In a preamble \cr ends a template only at align_state -1000000; "\span\iffalse{\fi" moves it back without a token in the template.
\tpl{a left brace skipped in the u template}{\iffalse{\span\iffalse}\fi#\fi}
\tpl{a right brace}{\iffalse\span\iffalse{\fi}#\fi}
\tpl{a left brace and a tab in the v template}{#\iffalse{&\span\iffalse}\fi\fi}
\tpl{a left brace and a cr}{#\iffalse{\cr\span\iffalse}\fi\fi}
\tpl{left, right, tab}{#\iffalse{}{&\span\iffalse}\fi\fi}
\noindent\par\message{^^Jright, tab:}
\setbox0\vbox{\halign{#\iffalse}&\span\iffalse{\fi\fi\cr x\cr}}\showbox0
\noindent\par\message{^^Jright, left, tab:}
\setbox0\vbox{\halign{#\iffalse}{&\span\iffalse}\fi\span\iffalse{\fi\fi\cr x\cr}}\showbox0
\noindent\par\message{^^Jright, left, cr:}
\setbox0\vbox{\halign{#\iffalse}{\cr\span\iffalse}\fi\span\iffalse{\fi\fi\cr x\cr}}\showbox0
\tpl{the end of the v template is skipped}{#\iffalse}

% Arguments in the templates.
\tpl{an undelimited argument across the u template's end}{\um{#}}
\tpl{a delimited argument across the u template's end}{\dm #.}
\tpl{a group in it, starting in the template}{\dm{#}.}
\tpl{an argument reading the end of the v template}{#\um}
\tpl{a delimited argument reading it}{#\dm a}
\tpl{in a group reading it}{#\um{a\span\iffalse}\fi}
\tpl{in a nested group reading it}{#\um{a{b}{c\span\iffalse}}\fi}
\tpl{a macro in the u template taking the cell}{\pass#}
\tpl{and a tab in its v part}{\um{a#b}&\um{c#d}}
\noindent\par
\setbox0\vbox{\halign{\um{#}&\dm #.\cr\omit x&\omit y\cr}}\showbox0
\setbox0\vbox{\halign{\um{#}\cr\noalign{\um{in noalign}\global\let\nn=\relax}x\cr}}\showbox0
\everycr{\um{in everycr}\noalign{\kern1pt}}
\setbox0\vbox{\halign{#\cr x\cr y\cr}}\showbox0
\everycr{}
\noindent\par

% Preamble: align_state is -1000000, and macros are expanded there only after \span.
\setbox0\vbox{\halign{\span\pass{a#}\span\um{#}\cr x\cr}}\showbox0
\setbox0\vbox{\halign{\span\dm{a}.#\cr x\cr}}\showbox0
\def\pr#1{#1&#1}
\setbox0\vbox{\halign{\span\pr{#}\cr x&y\cr}}\showbox0
\def\pp{\iffalse&\fi}
\setbox0\vbox{\halign{\span\pp#\cr x\cr}}\showbox0
\def\pq{\dm a&b.}
\setbox0\vbox{\halign{\span\pq#\cr x\cr}}\showbox0
\def\ps{\def\pt{a&b}}
\setbox0\vbox{\halign{\span\ps#\pt\cr x&y\cr}}\showbox0
\message{^^Jat the end}
\end
